\section{Metodologia}

\subsection{Contrato experimental e estrutura geral}

A comparação adota um único \textit{host} físico, um único artefato de
modelo verificado por \textit{hash} SHA-256 e um único contrato de
requisições HTTP, mantendo constante tudo que não seja a plataforma de
\textit{serving} ou o eixo experimental em estudo. Os dois pipelines
(FastAPI com ONNX Runtime e BentoML com \textit{adaptive batching})
expõem o mesmo \textit{endpoint} \texttt{POST /predict} que recebe um
\textit{JSON} com a chave \texttt{features} contendo um vetor bidimensional
de \texttt{float32} de forma $[N, 28]$ e retorna um vetor de probabilidades
$[N, 2]$. Ambos verificam o \textit{hash} do artefato ONNX no \textit{startup}
e expõem um \textit{endpoint} \texttt{/version} que reporta os
hiperparâmetros de paralelismo efetivamente aplicados; esse \textit{payload}
é registrado por célula da varredura para evitar atribuição errônea de
resultados.

O experimento se divide em três camadas de medição com finalidades
complementares: a primeira isola o custo intrínseco da inferência, a segunda
isola o sobrecusto da plataforma sem contêiner, e a terceira mede a entrega
ponta a ponta sob o contrato real de \textit{deployment}.

\subsection{Decomposição em três camadas}

\begin{table}[H]
\centering
\caption{As três camadas de medição.}
\label{tab:layers}
\begin{tabularx}{\textwidth}{l X X}
\toprule
Camada & O que isola & Como mede \\
\midrule
L1 & Custo intrínseco da inferência ONNX & Laço apertado em Python, 10\,000 iterações sobre \texttt{session.run}, sem HTTP, com \textit{cache} quente e \textit{ndarray} pré-construído. \\
L2 & Sobrecusto do \textit{framework} e do HTTP no \textit{host} & Servidor executado nativamente no \textit{host} com \texttt{taskset -c 0,1}; Vegeta hospedado nas \textit{cores} 2,3 ataca diretamente o \textit{loopback}. \\
L3 & \textit{Deployment} completo conteinerizado & Servidor em Docker com \texttt{--cpuset-cpus=0,1 --cpus=2 --memory=2g}; Vegeta nativo nas \textit{cores} 2,3. \\
\bottomrule
\end{tabularx}
\end{table}

A atribuição de sobrecusto é direta: L1 fornece o piso da inferência, $L2 -
L1$ representa o sobrecusto do \textit{framework}, do \textit{parsing}
\textit{JSON}, do despacho assíncrono e do tráfego TCP em \textit{loopback};
$L3 - L2$ representa o sobrecusto adicional do isolamento do contêiner e do
mapeamento de portas. Como o gerador de carga é o mesmo (Vegeta) nas três
camadas, as diferenças entre camadas não são contaminadas pela
implementação do cliente.

\subsection{Varredura matricial e variantes do BentoML}

A camada L3, que produz o resultado principal, percorre uma matriz
fatorial de configurações de paralelismo e intensidade de carga.

\begin{table}[H]
\centering
\caption{Eixos da varredura matricial L3.}
\label{tab:matrix}
\begin{tabularx}{\textwidth}{l X r}
\toprule
Eixo & Valores & Cardinalidade \\
\midrule
Variante & \texttt{fastapi}, \texttt{bentoml\_lat5}, \texttt{bentoml\_lat50}, \texttt{bentoml\_lat250} & 4 \\
(\textit{workers}, \textit{threads}) & $(1,1), (1,2), (2,1), (2,2)$ & 4 \\
RPS alvo & $\{50,\,150,\,350,\,600,\,1000\}$ com aborto precoce na saturação & até 5 \\
Repetições & 3 réplicas independentes & 3 \\
\bottomrule
\end{tabularx}
\end{table}

O eixo de variantes do BentoML promove o parâmetro \texttt{max\_latency\_ms}
de configuração fixa a eixo experimental explícito, depois que a leitura do
código-fonte do \textit{dispatcher} \cite{bentoml_dispatcher} mostrou que o
valor inicialmente calibrado (5\,ms) é um prazo rígido inferior ao piso de
\textit{serving} medido em L2. As variantes \texttt{lat50} e \texttt{lat250}
adicionam folga de uma e duas ordens de grandeza, respectivamente, permitindo
separar comportamento arquitetural de artefato de calibração. O parâmetro
\texttt{max\_batch\_size} é mantido em 16 para todas as variantes a fim de
isolar o efeito do prazo.

A varredura completa contabilizou 219 células efetivamente executadas, de
um máximo nominal de 240, com aborto precoce em valores de RPS além da
saturação observada. Cada célula registra um \textit{JSON} de versão, um
\textit{JSON} de resumo do Vegeta e o \textit{result.bin} bruto contendo a
latência e o código HTTP de cada requisição individual.

\subsection{Geração de carga em modelo aberto}

Vegeta \cite{vegeta2024} é configurado em modo de modelo aberto a taxa fixa,
com \textit{timeout} por requisição de 3 segundos, formato de saída
binário, e fonte de alvos pré-renderizada a partir de 10\,000 linhas
amostradas do conjunto de teste. Cada célula recebe 100 requisições de
aquecimento antes da medição, removendo a influência do regime transitório
de \textit{startup} do servidor (em especial a fase \texttt{train\_optimizer}
do \textit{dispatcher} do BentoML) e da inicialização tardia das estruturas
internas do \textit{tokenizer} e do \textit{thread pool}. A duração nominal
do ataque é de 45 segundos. Para reduzir efeitos de variabilidade entre
\textit{containers}, é introduzido um intervalo de quietude de 3 segundos
entre células consecutivas.

\subsection{Protocolo de pinagem de CPU}

\begin{table}[H]
\centering
\caption{Pinagem de CPU por camada.}
\label{tab:pinning}
\begin{tabularx}{\textwidth}{l X X}
\toprule
Camada & Servidor & Vegeta \\
\midrule
L1 & não se aplica & não se aplica \\
L2 & \texttt{taskset -c 0,1} (host) & nativo, \texttt{taskset -c 2,3} \\
L3 & Docker \texttt{--cpuset-cpus=0,1 --cpus=2} & nativo, \texttt{taskset -c 2,3} \\
\bottomrule
\end{tabularx}
\end{table}

A disjunção entre as \textit{cores} do servidor e as do gerador de carga é
necessária para impedir contaminação cruzada entre a carga injetada e a
carga inferida. As demais \textit{cores} (4 a 11 no \textit{host} de
referência) permanecem disponíveis para o \textit{kernel} Linux atender
demais processos, sem efeito mensurável sobre a medição.

\subsection{Amostragem de utilização de CPU no servidor}

Um subconjunto de 20 células consideradas críticas para o argumento do
mecanismo é re-executado com amostragem direta de utilização de CPU do
\textit{container} do servidor, via \texttt{docker stats --no-stream}
acionado em laço com período efetivo próximo de um segundo. Cada célula
desse subconjunto, repetida três vezes por 45 segundos, produz uma série
temporal por execução (cerca de 22 amostras) e seus resumos (média,
percentil 50, percentil 95 e máximo). O subconjunto cobre todas as
variantes em configurações próximas à saturação, totalizando 60 execuções
adicionais; o conjunto principal de 219 células é preservado intacto em
diretório separado.

\subsection{Procedimentos estatísticos e métrica de comparação}

As métricas reportadas por célula são as medianas entre as três réplicas.
Acompanhando cada percentil reportado, é computado o intervalo de
confiança bootstrap percentílico a 95\% sobre as três réplicas, com 1.000
reamostragens \cite{efron1979}, e o coeficiente de variação
$\text{CoV} = \sigma/\mu$, usado como indicador de estabilidade entre
execuções. Os intervalos resultantes são largos em termos absolutos a
$n = 3$, e as comparações entre pipelines apoiam-se na magnitude da
diferença observada, frequentemente de ordens de grandeza no P99, mais do
que na estreitez dos intervalos. Comparações pareadas adicionais entre
variantes empregam o teste de Mann-Whitney U e o tamanho de efeito de
Cliff's $\delta$, ambos não-paramétricos e adequados a amostras pequenas.

Para comparar de forma justa um sistema que descarta requisições por prazo
expirado (BentoML, HTTP 503) contra um sistema que as enfileira como
latência adicional (FastAPI, código 200 com latência elevada), a métrica
principal adotada é a vazão útil sob SLA, calculada como
$\text{goodput@SLA} = N_{\text{ok}}^{<\text{SLA}} / N_{\text{oferecidas}}$,
onde $N_{\text{ok}}^{<\text{SLA}}$ é o número de respostas 2xx com latência
menor ou igual ao limiar (50 ou 100\,ms). Os dois modos de falha (503
imediato e 200 atrasado) são tratados de forma simétrica como falhas
operacionais idênticas do ponto de vista do consumidor.

\subsection{Ameaças à validade}

Quatro ameaças à validade interna foram identificadas e tratadas. Primeira,
o agrupamento de conexões HTTP/1.1 com \textit{keep-alive} do Vegeta tende
a colapsar para uma única conexão TCP sob latência mediana submilissegundo,
o que efetivamente fixa a carga a um único \textit{worker} do FastAPI sob
balanceamento por conexão. Esse efeito foi confirmado empiricamente nas
células do subconjunto de mecanismo e é discutido como limitação na
Seção~7. Segunda, o \textit{dispatcher} do BentoML possui uma fase de
\textit{warmup} de cerca de 30 segundos durante a qual o estimador de
tamanho ótimo de lote estabiliza; o aquecimento de 100 requisições foi
calibrado para drenar essa fase. Terceira, a migração entre \textit{Vegeta}
em \textit{Docker} (Windows) e \textit{Vegeta} nativo (WSL2 Linux) altera o
caminho do \textit{loopback}; os resultados aqui reportados usam o caminho
nativo, e o conjunto Windows é preservado em diretório separado para
verificação cruzada. Quarta, o trabalho restringe-se a um \textit{host}
físico; uma extensão multi-\textit{host} com balanceador de carga real é
apontada como trabalho futuro.

